\section{Demokratisches Gegengewicht} \label{sec:demokratisches-gegengewicht}
Durch das vorgeschlagene Wahlverfahren bekommt die Regierung mehr Macht.
Dafür muss an anderer Stelle ein demokratischer Ausgleich geschaffen werden, damit die Gefahr des Festsetzens an der Macht wieder verringert wird.

Die Regierung hat mehr Macht bekommen, damit sie ihren Auftrag, zu gestalten, auch in herausfordernden Zeiten, wo viel entschieden werden muss, erfüllen kann.
Umso wichtiger ist es, dass es eine "demokratische Notbremse" gibt, mit der man eine schlechte Regierung auch vorzeitig wieder loswerden kann, wenn diese ihre Macht missbraucht oder durch Inkompetenz
verschwendet.
Es muss Mechanismen geben, eine vorzeitige Ablösung von Einzelpersonen oder der ganzen Regierung einzuleiten, wenn dies geboten ist.

Die Herausforderung ist, das Machtgleichgewicht so zu gestalten, dass eine Regierung einerseits genug Macht hat, um auch unpopuläre Entscheidungen durchzusetzen, und andererseits noch gestoppt werden kann,
wenn sie zu viel Schaden anrichtet.

Im Folgenden sollen daher verschiedene Fälle diskutiert werden:

\subsection[]{Einzelne Parlamentarier verstoßen gegen das Gesetz}

Einzelne Parlamentarier können bereits jetzt nach Art 46 GG strafrechtlich verfolgt werden, wenn der Bundestag dem zustimmt.
Darüberhinaus sollte der Bundestag das Recht haben, Abgeordnete für die Dauer des Verfahrens zu suspendieren und bei
nachgewiesenen schweren Vergehen das Mandat sogar rückwirkend zu entziehen, so dass der Abgeordnete alle Privilegien
verliert, die mit einem Bundestagsmandat verbunden sind.

Schwere Vergehen in diesem Sinne sind Vergehen
\begin{itemize}
  \item die sich vorsätzlich gegen die Bundesrepublik Deutschland richten, wie z. B. Spionage für eine fremde Macht, oder
  \item Vergehen, die nach allgemeinem Rechtsempfinden ein Volksvertreteramt ausschließen, wie z. B. Gewaltverbrechen oder Sexualstraftaten.
 \end{itemize}

99\% der Bevölkerung würden sich wohl nicht von einem Mörder oder Kinderschänder vertreten lassen wollen, auch wenn er anderweitig noch so brilliant ist.
Bei Steuerhinterziehung und Trunkenheitsfahrten hingegen könnten die Meinungen vielleicht geteilt sein.

Der Bundestag hat das Recht bei Nachweis eines schweren Vergehens, das Mandat ab sofort zu entziehen oder sogar rückwirken, wenn der Strafbestand schon die ganze
Legislaturperiode bestanden hat.
Der betroffene Abgeordnete hat das Recht, bei dem Bundesverfassungsgericht gegen diese Maßnahme zu klagen.
Gibt es umgekehrt im Bundestag keine Mehrheit für die Suspendierung oder den Mandatsentzug, können einzelne Abgeordnete beim Bundesverfassungsgericht darauf klagen, und das
Bundesverfassungsgericht kann der Klage entsprechen oder nicht.


 {abloesung_einzelner_parlamentarier}